\documentclass[spanish,letterpaper,oneside]{article}

\def\documenttitle {¿Cómo Soy Como Maestro? Silueta Docente e Identidad Reflexiva}
\def\documentsubtitle {Fundamentos para la enseñanza y el aprendizaje I · Cohorte SD26}
\def\documentsubject {Actividad 2 · Sesión 2}
\def\documentauthor {Mtro. Martín Jonathan de la Cruz Muñoz}
\def\coursename {Fundamentos para la enseñanza y el aprendizaje I}
\def\coursecode {MAEM-SD26}

\def\universityname {Instituto de Investigación, Innovación y Estudios de Posgrado para la Educación del Estado de Nuevo León}
\def\universityfaculty {Maestría en Aprendizaje y Enseñanza de las Matemáticas}
\def\universitydepartment {Fundamentos para la enseñanza y el aprendizaje I}
\def\universitydepartmentimage {img/iiiepe/marca-iiiepe.png}
\def\universitydepartmentimagecfg {height=1.75cm}
\def\universitylocation {Monterrey, Nuevo León}
\def\studentfullname {Mtro. Martín Jonathan de la Cruz Muñoz}
\def\studentid {00874}
\def\studentemail {martin.delacruz@campus.iiiepe.edu.mx}
\def\programperiod {2do periodo, septiembre--diciembre de 2026 (SD26)}
\def\programmodality {Modalidad presencial}
\def\programshort {MAEM}

\def\authortable {
  \begin{tabular}{ll}
    \textbf{Alumno:} & \studentfullname \\
    Matrícula: & \studentid \\
    Correo institucional: & \studentemail \\
    Unidad: & \coursename \\
    Programa: & \programshort \\
    Periodo: & \programperiod \\
    & \\
    \multicolumn{2}{l}{Fecha de realización: 19 de noviembre de 2026} \\
    \multicolumn{2}{l}{\universitylocation}
  \end{tabular}
}

\input{template}

\usepackage{helvet}
\renewcommand{\familydefault}{\sfdefault}
\usepackage{array}
\usepackage{longtable}
\usepackage{pdflscape}
\usepackage{tikz}
\usepackage{pgfplots}
\pgfplotsset{compat=1.18}
\usepackage[most]{tcolorbox}
\usetikzlibrary{arrows.meta,positioning,calc,fit,backgrounds,shapes.geometric}

\definecolor{iiiepeNavy}{HTML}{163B5C}
\definecolor{iiiepeBlue}{HTML}{1F6F9F}
\definecolor{iiiepeAqua}{HTML}{20B7A8}
\definecolor{iiiepeOrange}{HTML}{F28C28}
\definecolor{iiiepePaper}{HTML}{F6F8F7}
\definecolor{iiiepeInk}{HTML}{1E2A32}
\definecolor{iiiepePale}{HTML}{EEF5F8}
\definecolor{iiiepeMint}{HTML}{E7F6F4}
\definecolor{iiiepeWarm}{HTML}{FFF3E6}
\def\portraittitlecolor {iiiepeNavy}
\setcitestyle{authoryear,open={(},close={)}}
\makeatletter
\renewcommand{\@biblabel}[1]{}
\makeatother
\setlength{\bibhang}{1.27cm}
\setlength{\bibsep}{6pt}
\AtBeginDocument{\renewcommand{\refname}{Referencias}}

\newcommand{\insertcoverwatermark}{%
  \AddToShipoutPictureBG*{%
    \begin{tikzpicture}[remember picture,overlay]
      \node[opacity=0.04,inner sep=0pt] at (current page.center) {%
        \includegraphics[width=0.65\paperwidth]{img/iiiepe/marca-iiiepe.png}%
      };
    \end{tikzpicture}%
  }%
}

\setlength{\parskip}{5pt}\setlength{\parindent}{0pt}
\begin{document}\insertcoverwatermark\templatePortrait\templatePagecfg\templateFinalcfg\fancyfoot{}\fancyfoot[L]{\scriptsize IIIEPE · MAEM · SD26}\fancyfoot[C]{\scriptsize\thepage}\fancyfoot[R]{\scriptsize Actividad 2 · Identidad Docente}\setstretch{1.3}\raggedbottom
\section{Introducción y postura central sobre la identidad docente}
Reconocer cómo enseñamos implica mirar más allá de las técnicas empleadas en clase para interrogar las decisiones didácticas, los criterios éticos y las creencias epistemológicas que orientan nuestra labor cotidiana. A partir de la reflexión crítica sobre el interrogante \textit{¿cómo soy como maestro?}, sostengo la tesis de que la identidad docente en matemáticas se construye en la tensión productiva entre el rigor técnico-disciplinario heredado de la ingeniería y una pedagogía humanista, dialógica y socialmente responsable \citep{medina2009didactica,unesco2015replantear}.

El análisis realizado me lleva a sostener que no basta con identificar virtudes o carencias de manera abstracta. Por ello, este documento responde primero de forma argumentada a las cinco preguntas de reflexión inicial, despliega a continuación una \textbf{silueta docente} gráficamente estructurada que sintetiza centro identitario, fortalezas y áreas de oportunidad situadas en episodios concretos de aula, así como necesidades de formación bibliográfica justificadas, y culmina con líneas de transferencia práctica e implicaciones formativas para mi tránsito en la MAEM.

\section{Reflexión inicial: cinco preguntas sobre la práctica de enseñanza}

\subsection*{1. ¿Cómo describes tu forma de enseñar?}
Describo mi forma de enseñar como una transición consciente y gradual: provengo de una cultura de organización técnica, industrial y de ingeniería —donde priman el orden procedimental, la precisión en magnitudes y la solución eficiente de fallas—, pero reconozco que ese rigor pierde sentido educativo si se vuelve vertical o inaccesible para el grupo. Por consiguiente, procuro estructurar la clase combinando claridad operativa con preguntas abiertas que inviten a comprender el significado físico y lógico de cada expresión matemática \citep{pimienta2012estrategias}.

\subsection*{2. ¿Qué estrategias utilizas con mayor frecuencia?}
Con mayor frecuencia empleo la resolución contextualizada de problemas de taller o ingeniería (balances sencillos, relaciones proporcionales, ecuaciones de calibración y lectura de gráficas experimentales), el trabajo guiado paso a paso con verificación de unidades, el contraste entre representaciones tabular, gráfica y algebraica, así como el intercambio en binas o equipos pequeños para comparar rutas de solución \citep{star2015algebra}. Sin embargo, advierto la necesidad de sistematizar estrategias de argumentación escrita y evaluación formativa con criterios más explícitos \citep{eef2021feedback}.

\subsection*{3. ¿Cuáles consideras que son tus principales fortalezas como docente?}
Considero que mis fortalezas radican en tres ámbitos: primero, el dominio aplicado de las matemáticas en entornos de mantenimiento, mecatrónica e ingeniería, lo cual me permite diseñar ejemplos auténticos y responder con solvencia cuando el alumnado pregunta para qué sirve un contenido; segundo, la disposición para organizar secuencias estructuradas y materiales visuales claros; y tercero, la convicción ética de tratar con respeto y dignidad a cada estudiante independientemente de sus rezagos previos.

\subsection*{4. ¿Qué aspectos de tu práctica te gustaría mejorar?}
Me interesa mejorar principalmente tres dimensiones: en primer lugar, reducir la tentación de intervenir demasiado rápido resolviendo el paso crítico en el pizarrón cuando un equipo se detiene, cediendo mayor protagonismo cognitivo al estudiantado \citep{sawyer2008optimizing}; en segundo lugar, diseñar instrumentos de retroalimentación descriptiva que atiendan la diversidad de ritmos sin limitarse a calificar aciertos finales \citep{bostrom2023formative}; y en tercer lugar, documentar sistemáticamente evidencias de aula para tomar decisiones basadas en indagación pedagógica \citep{restrepo2004investigacion}.

\subsection*{5. ¿Cómo tomas decisiones durante tu enseñanza?}
Tradicionalmente tomaba decisiones en el aula guiado por la lógica del temario y la intuición acumulada ante las dudas inmediatas del grupo: si percibía silencio o error generalizado, repetía el procedimiento con otro ejemplo similar. No obstante, este análisis propio me lleva a replantear ese criterio: una decisión didáctica fundada requiere interpretar qué obstáculo conceptual originó la confusión, ofrecer un contraejemplo o andamiaje pertinente y revisar posteriormente si el ajuste amplió realmente la comprensión colectiva \citep{eef2017mathematics}.

\section{Silueta docente: identidad, prácticas situadas y horizonte formativo}
La Figura~\ref{fig:act2-silueta} representa de manera integrada los tres ejes constitutivos de la identidad profesional docente: en el centro, los rasgos que definen mi identidad profesional; en el lado izquierdo, fortalezas y áreas de oportunidad explicadas en situaciones reales de enseñanza técnica y universitaria; y en el lado derecho, las necesidades de formación y fuentes bibliográficas justificadas para fortalecer mi docencia.

\begin{figure}[htbp]
\centering
\begin{tikzpicture}[font=\scriptsize, scale=0.88, transform shape]
  % Zona central: Silueta humana limpia y amplia sin traslapes laterales
  % Cabeza y cuello
  \draw[fill=iiiepePale,draw=iiiepeNavy,very thick] (0,3.15) circle (0.72cm);
  \node[align=center,font=\scriptsize\bfseries,iiiepeNavy] at (0,3.15) {RIGOR\\Y ÉTICA};

  % Torso y brazos estilizados contenidos dentro de x in [-1.15, 1.15]
  \draw[fill=iiiepeBlue!10,draw=iiiepeNavy,very thick,rounded corners=10pt]
    (-0.32,2.42) -- (0.32,2.42) -- (1.12,2.10) -- (1.02,0.55) -- (0.78,0.55) -- (0.72,1.65)
    -- (0.66,-1.95) -- (0.14,-1.95) -- (0.06,-0.25) -- (-0.06,-0.25) -- (-0.14,-1.95) -- (-0.66,-1.95)
    -- (-0.72,1.65) -- (-0.78,0.55) -- (-1.02,0.55) -- (-1.12,2.10) -- cycle;

  % Etiquetas centrales posicionadas relativamente (cero traslape)
  \node[draw=iiiepeNavy,fill=white,rounded corners=3pt,thick,text width=1.95cm,align=center,inner sep=3pt,anchor=north] (c1) at (0,2.28)
    {\textbf{\color{iiiepeNavy}CENTRO}\\Identidad técnica y humanista};
  \node[draw=iiiepeAqua!80!black,fill=iiiepeMint,rounded corners=3pt,thick,text width=1.85cm,align=center,inner sep=3pt,below=0.22cm of c1] (c2)
    {\textbf{\color{iiiepeNavy}RASGOS CLAVE}\\Rigor · Diálogo\\Dignidad · Equidad};

  % Columna izquierda: x <= -1.65 y posicionamiento vertical relativo
  \node[draw=iiiepeNavy,fill=iiiepePale,rounded corners=4pt,thick,text width=5.65cm,align=left,inner sep=6pt,anchor=north east] (l1) at (-1.65,3.85) {%
    \textbf{\color{iiiepeNavy}\small LADO IZQUIERDO: FORTALEZAS SITUADAS}\\[2pt]
    \textbf{1. Modelación técnica aplicada (CONALEP/CECyTE):} Al estudiar relaciones lineales con sensores o costos de mantenimiento, traduzco fenómenos físicos en modelos matemáticos verificables.\\[2pt]
    \textbf{2. Claridad estructural (Ingeniería):} Organizo problemas complejos en etapas de análisis dimensional, hipótesis y comprobación crítica.\\[2pt]
    \textbf{3. Trato digno e inclusivo:} Genero confianza en el pizarrón validando las preguntas genuinas sin ridiculizar bloqueos previos.
  };

  \node[draw=iiiepeOrange!90!black,fill=iiiepeWarm,rounded corners=4pt,thick,text width=5.65cm,align=left,inner sep=6pt,below=0.28cm of l1.south west,anchor=north west] (l2) {%
    \textbf{\color{iiiepeOrange!85!black}\small ÁREAS DE OPORTUNIDAD EN EL AULA}\\[2pt]
    \textbf{1. Prisa por cerrar el algoritmo:} Ante falta de tiempo de bloque, he terminado mostrando el despeje antes de que los equipos conjeturen su estrategia.\\[2pt]
    \textbf{2. Evaluación formativa sistemática:} Paso con facilidad de ejercicios guiados a exámenes escritos sin matrices previas de argumentación.\\[2pt]
    \textbf{3. Heterogeneidad matemática:} Requiero andamiajes escalonados para alumnos que ingresan con fragilidad aritmética y algebraica.
  };

  % Columna derecha: x >= +1.65 y posicionamiento vertical relativo
  \node[draw=iiiepeBlue,fill=iiiepePale,rounded corners=4pt,thick,text width=5.65cm,align=left,inner sep=6pt,anchor=north west] (r1) at (1.65,3.85) {%
    \textbf{\color{iiiepeBlue}\small LADO DERECHO: NECESIDADES DE FORMACIÓN}\\[2pt]
    \textbf{1. Didáctica para el aprendizaje profundo:} Diseñar secuencias basadas en indagación y andamiaje progresivo \citep{sawyer2008optimizing}.\\[2pt]
    \textbf{2. Retroalimentación que haga pensar:} Dominar devoluciones descriptivas y criterios compartidos \citep{eef2021feedback}.\\[2pt]
    \textbf{3. Investigación-acción e inclusión:} Documentar evidencia de clase para adaptar la enseñanza a contextos sociofamiliares diversos \citep{restrepo2004investigacion,bronfenbrenner2006bioecologico}.
  };

  \node[draw=iiiepeAqua!80!black,fill=iiiepeMint,rounded corners=4pt,thick,text width=5.65cm,align=left,inner sep=6pt,below=0.28cm of r1.south west,anchor=north west] (r2) {%
    \textbf{\color{iiiepeAqua!75!black}\small BIBLIOGRAFÍA PRIORITARIA JUSTIFICADA}\\[2pt]
    \textbf{· \citeauthor{sawyer2008optimizing} (2008):} Comprender principios de ciencias del aprendizaje frente a la transmisión.\\[2pt]
    \textbf{· \citeauthor{meirieu2020escuela} (2020):} Reconstruir el sentido democrático, colectivo y emancipador del aula.\\[2pt]
    \textbf{· \citeauthor{restrepo2004investigacion} (2004):} Convertir la propia clase en objeto de indagación rigurosa.\\[2pt]
    \textbf{· \citeauthor{bronfenbrenner2006bioecologico} (2006):} Situar al estudiante dentro de sus ecosistemas reales.
  };

  % Conectores limpios sin cruzar el texto
  \draw[->,thick,iiiepeNavy] (l1.east) -- (-0.95,1.95);
  \draw[->,thick,iiiepeOrange!90!black] (l2.east) -- (-0.68,-0.80);
  \draw[->,thick,iiiepeBlue] (1.05,1.95) -- (r1.west);
  \draw[->,thick,iiiepeAqua!80!black] (0.68,-0.80) -- (r2.west);
\end{tikzpicture}
\caption{Silueta docente: articulación entre identidad profesional, análisis situado de la práctica y ruta de formación.}
\label{fig:act2-silueta}
\end{figure}

\section{Desarrollo analítico de los tres componentes de la silueta}

\subsection*{Centro de la silueta: autodefinición y elementos característicos}
Me defino como un docente en reconstrucción reflexiva cuya base técnica e ingenieril aporta estructura lógica, aprecio por los datos verificables y vínculo con el entorno productivo \citep{conalep2025mecatronica}, mientras su proceso formativo en la MAEM orienta esos recursos hacia la dignidad del estudiante, el diálogo respetuoso y el aprendizaje comprensivo. Tres elementos caracterizan mi práctica: \textbf{rigor} (cuidar consistencia lógica, unidades y validez de los supuestos), \textbf{diálogo} (validar la palabra del alumnado para descubrir cómo piensa) y \textbf{dignidad} (sostener altas expectativas acompañadas de apoyos accesibles, nunca desde el castigo).

\subsection*{Lado izquierdo: fortalezas y áreas de oportunidad en situaciones reales}
Para evitar enumeraciones descontextualizadas, sintetizo la manifestación concreta de ambos polos en el cuadro comparativo siguiente:

{\small\setstretch{1.15}\renewcommand{\arraystretch}{1.35}
\begin{longtable}{|>{\raggedright\arraybackslash}p{0.20\linewidth}|>{\raggedright\arraybackslash}p{0.39\linewidth}|>{\raggedright\arraybackslash}p{0.35\linewidth}|}\hline
\rowcolor{iiiepePale} \textbf{Dimensión} & \textbf{Manifestación situada en clase} & \textbf{Interpretación pedagógica} \\\hline
\endfirsthead
\hline
\rowcolor{iiiepePale} \textbf{Dimensión} & \textbf{Manifestación situada en clase} & \textbf{Interpretación pedagógica} \\\hline
\endhead
Fortaleza: Modelación técnica situada & En grupos de bachillerato técnico en mantenimiento y mecatrónica, al abordar variación lineal con lecturas de calibración y costos de servicio, logro que estudiantes habitualmente distantes del álgebra reconozcan el sentido práctico de la pendiente y las intersecciones. & Vincula abstracción matemática con referentes físicos tangibles, activando motivación intrínseca y pertinencia laboral \citep{star2015algebra}. \\\hline
Fortaleza: Acompañamiento respetuoso & En cursos de ingeniería, frente a errores de despeje o conversión de unidades en el pizarrón, detengo la burla o el silencio punitivo pidiendo al grupo rastrear en qué paso cambió el significado físico de la expresión. & Convierte el error en objeto de análisis colectivo, reduciendo la ansiedad matemática \citep{schunk2012teorias}. \\\hline
Área de oportunidad: Toma de control prematura & Al observar que el cierre de la sesión se aproxima y varios equipos dudan ante un sistema de ecuaciones, tiendo a resolver el ejemplo en el pizarrón para «no atrasarme» con la planeación. & Esto implica priorizar la cobertura temática por encima de la consolidación autónoma del razonamiento estudiantil \citep{sawyer2008optimizing}. \\\hline
Área de oportunidad: Retroalimentación tardía & En ocasiones concentro la verificación de avances en pruebas escritas parciales, recibiendo señales de incomprensión cuando ya concluyó la unidad didáctica. & Demanda transitar hacia evidencias cortas frecuentes y listas de cotejo formativas durante el proceso \citep{eef2021feedback}. \\\hline
\end{longtable}}

\subsection*{Lado derecho: necesidades de formación y bibliografía justificada}
Reconocer dichas tensiones fundamenta una agenda sistemática de formación profesional en la MAEM. Además de evitar lecturas aisladas, cada fuente elegida responde directamente a un obstáculo observado en clase; en el cuadro siguiente explico qué tensión atiende cada necesidad formativa y qué aportación práctica espero consolidar en mis grupos:

{\small\setstretch{1.15}\renewcommand{\arraystretch}{1.32}
\begin{longtable}{|>{\raggedright\arraybackslash}p{0.23\linewidth}|>{\raggedright\arraybackslash}p{0.31\linewidth}|>{\raggedright\arraybackslash}p{0.40\linewidth}|}\hline
\rowcolor{iiiepePale} \textbf{Necesidad / Autor clave} & \textbf{Tensión docente atendida} & \textbf{Justificación y meta práctica en aula} \\\hline
\endfirsthead
\hline
\rowcolor{iiiepePale} \textbf{Necesidad / Autor clave} & \textbf{Tensión docente atendida} & \textbf{Justificación y meta práctica en aula} \\\hline
\endhead
Andamiaje comprensivo · \citet{sawyer2008optimizing} & Tendencia a resolver el despeje en pizarrón cuando los equipos se detienen. & Aprender a graduar pistas y preguntas detonadoras que retengan el protagonismo cognitivo del estudiantado. \\\hline
Pedagogía cooperativa y derecho a aprender · \citet{meirieu2020escuela} & Desigualdad de saberes previos en bachillerato técnico e ingeniería. & Diseñar dinámicas de diálogo colectivo (\textit{yo + nosotros}) y rodeos didácticos (\textit{algo más}) en lugar de repetir recetas. \\\hline
Evaluación formativa oportuna · \citet{eef2021feedback,pimienta2008evaluacion} & Concentrar la verificación en exámenes parciales tardíos sin devolución procesual. & Instalar listas de valoración claras, retroalimentación descriptiva temprana y segundas versiones corregidas. \\\hline
Indagación docente sistemática · \citet{restrepo2004investigacion} & Tomar decisiones curriculares basadas solo en intuición inmediata ante el grupo. & Documentar cuadernos, errores frecuentes y diarios de clase como ciclos formales de mejora continua. \\\hline
Contexto ecológico integral · \citet{bronfenbrenner2006bioecologico} & Subestimar jornadas laborales, transporte y condiciones familiares en la planeación. & Diseñar guías accesibles y procesos proximales sostenidos que no conviertan el contexto en filtro de exclusión \citep{bostrom2023formative}. \\\hline
\end{longtable}}

\section{Transferencia profesional y proyecciones prácticas}
En el contexto laboral de educación media superior y superior tecnológica, el impacto de este análisis reflexivo consiste en convertir la autoobservación en decisiones concretas de planeación: reservar los últimos diez minutos de cada bloque para que al menos dos equipos argumenten sus estrategias sin intervención anticipada del docente, explicitar criterios de evaluación antes de iniciar cada problema abierto e incorporar cuestionarios diagnósticos breves sobre comprensión conceptual además del cálculo numérico \citep{bostrom2023formative}.

\clearpage
\section{Conclusiones y síntesis argumentativa}
En conclusión, responder con honestidad a la pregunta \textit{¿cómo soy como maestro?} revela que la identidad profesional no es un rasgo fijo ni terminado, sino una construcción ética e intelectual continua. La silueta analizada pone de manifiesto que mi trayectoria técnica constituye un activo valioso siempre que se articule con la escucha pedagógica, la cesión gradual de autonomía y la evaluación formativa.

Por tanto, el compromiso asumido en esta cohorte SD26 consiste en avanzar desde una enseñanza basada en la exposición ordenada hacia una mediación capaz de propiciar pensamiento matemático profundo, equitativo y transferible, documentando cada cambio con evidencia reflexiva a lo largo de los bloques siguientes.

\clearpage
\begin{thebibliography}{99}
\bibitem[Boström y Palm(2023)]{bostrom2023formative} Boström, E., y Palm, T. (2023). Formative assessment practice in mathematics: Teachers' perceptions of prerequisite conditions. \textit{Frontiers in Education}, 8, Artículo 1101192. \url{https://doi.org/10.3389/feduc.2023.1101192}
\bibitem[Bronfenbrenner y Morris(2006)]{bronfenbrenner2006bioecologico} Bronfenbrenner, U., y Morris, P. A. (2006). The bioecological model of human development. En R. M. Lerner (Ed.), \textit{Handbook of child psychology: Vol. 1. Theoretical models of human development} (6.ª ed., pp. 793--828). John Wiley \& Sons.
\bibitem[Collin y Quigley(2021)]{eef2021feedback} Collin, J., y Quigley, A. (2021). \textit{Teacher feedback to improve pupil learning}. Education Endowment Foundation.
\bibitem[CONALEP(2025)]{conalep2025mecatronica} Colegio Nacional de Educación Profesional Técnica. (2025). \textit{Perfil de egreso: Profesional Técnico-Bachiller en Mecatrónica. Modelo Académico 2025}. Secretaría Académica, CONALEP.
\bibitem[Henderson et~al.(2017)]{eef2017mathematics} Henderson, P., Hodgen, J., Foster, C., y Kuchemann, D. (2017). \textit{Improving mathematics in Key Stages 2 and 3}. Education Endowment Foundation.
\bibitem[Medina Rivilla y Salvador Mata(2009)]{medina2009didactica} Medina Rivilla, A., y Salvador Mata, F. (2009). \textit{Didáctica general} (2.ª ed.). Pearson Educación.
\bibitem[Meirieu(2020)]{meirieu2020escuela} Meirieu, P. (2020, 18 de abril). ``La escuela después''... ¿Con la pedagogía de antes? \textit{Movimiento Cooperativo de Escuela Popular}. \url{https://www.mcep.es/la-escuela-despues-con-la-pedagogia-de-antes-philippe-meirieu/}
\bibitem[Pimienta Prieto(2008)]{pimienta2008evaluacion} Pimienta Prieto, J. H. (2008). \textit{Evaluación de los aprendizajes: Un enfoque basado en competencias}. Pearson Educación.
\bibitem[Pimienta Prieto(2012)]{pimienta2012estrategias} Pimienta Prieto, J. H. (2012). \textit{Estrategias de enseñanza-aprendizaje: Docencia universitaria basada en competencias}. Pearson Educación.
\bibitem[Restrepo Gómez(2004)]{restrepo2004investigacion} Restrepo Gómez, B. (2004). La investigación-acción educativa y la construcción de saber pedagógico. \textit{Educación y Educadores}, (7), 45--55.
\bibitem[Sawyer(2008)]{sawyer2008optimizing} Sawyer, R. K. (2008). Optimizing learning: Implications of learning sciences research. En OECD, \textit{Innovating to learn, learning to innovate} (pp. 45--65). OECD Publishing.
\bibitem[Schunk(2012)]{schunk2012teorias} Schunk, D. H. (2012). \textit{Teorías del aprendizaje: Una perspectiva educativa} (6.ª ed.). Pearson Educación.
\bibitem[Star et~al.(2015)]{star2015algebra} Star, J. R., Caronongan, P., Foegen, A., Furgeson, J., Keating, B., Larson, M. R., Lyskawa, J., McCallum, W. G., Porath, J., y Zbiek, R. M. (2015). \textit{Teaching strategies for improving algebra knowledge in middle and high school students} (NCEE 2015-4010). Institute of Education Sciences, U.S. Department of Education.
\bibitem[UNESCO(2015)]{unesco2015replantear} UNESCO. (2015). \textit{Replantear la educación: ¿Hacia un bien común mundial?} Ediciones UNESCO.
\end{thebibliography}

\end{document}
